% Реферат: Моделювання турбулентності
% Оформлення: ДСТУ 3008:2015 (поля, шрифт, інтервал, підписи)
% Компіляція: xelatex turbulence_modeling_report.tex && xelatex turbulence_modeling_report.tex
\documentclass[14pt,a4paper,extarticle]{extarticle}
\usepackage{fontspec}
\usepackage{polyglossia}
\setdefaultlanguage{ukrainian}
\setmainfont{Times New Roman}[
  Ligatures=TeX,
  BoldFont = {Times New Roman Bold},
  ItalicFont = {Times New Roman Italic},
  BoldItalicFont = {Times New Roman Bold Italic},
]
\usepackage{amsmath, amssymb, amsfonts}
\usepackage{geometry}
\usepackage{graphicx}
\usepackage{hyperref}
\usepackage{booktabs}
\usepackage{array}
\usepackage{float}
\usepackage{caption}
\usepackage{setspace}
\usepackage{indentfirst}

% ДСТУ 3008:2015 — поля (ліве підшиття 30 мм)
\geometry{left=3cm, right=1cm, top=2cm, bottom=2cm}
\onehalfspacing
\setlength{\parindent}{1.25cm}
\setlength{\parskip}{0pt}

\hypersetup{colorlinks=true, linkcolor=black, urlcolor=black, citecolor=black}

\renewcommand{\contentsname}{Зміст}
\renewcommand{\figurename}{Рис.}
\renewcommand{\tablename}{Таблиця}
\DeclareCaptionLabelSeparator{ukrdash}{~---~}
\captionsetup{
  font=small,
  labelfont=bf,
  labelsep=endash,
  skip=6pt,
}
\captionsetup[table]{
  position=top,
  skip=0pt,
  font=normalsize,
  labelfont=normalfont,
  labelsep=ukrdash,
  justification=raggedright,
  singlelinecheck=false,
}

\newcommand{\Rey}{\mathrm{Re}}

\begin{document}

% -------------------------------------------------------------------
\begin{titlepage}
  \centering\thispagestyle{empty}
  Міністерство освіти і науки України\\[0.4em]
  Київський національний університет імені Тараса Шевченка\\[0.4em]
  Факультет комп'ютерних наук та кібернетики
  \vfill
  {\Large\bfseries РЕФЕРАТ}\\[0.6em]
  з дисципліни\\[0.2em]
  <<Чисельне моделювання динаміки систем>>\\[1.2em]
  {\large\bfseries Тема: Моделювання турбулентності:\\
  фундаментальні фізико-математичні аспекти,\\
  замикання рівнянь Нав'є--Стокса та чисельні методи}
  \vfill
  \begin{flushright}
    \begin{tabular}{@{}l@{}}
      Виконав:\\
      Кіщук Я.~Я.\\[1.2em]
    \end{tabular}
  \end{flushright}
  \vfill
  Київ~--- 2026
\end{titlepage}

% -------------------------------------------------------------------
\section*{Анотація}
\addcontentsline{toc}{section}{Анотація}

У рефераті систематизовано фізичні властивості турбулентного руху рідини,
теорію енергетичного каскаду та роль числа Рейнольдса. Проаналізовано
обчислювальні вимоги до прямого чисельного моделювання (DNS), процедуру
осереднення Рейнольдса та проблему замикання RANS-моделей. Наведено
класифікацію алгебраїчних, двопараметричних ($k$--$\varepsilon$,
$k$--$\omega$, SST) та LES-підходів; обговорено вплив схемної в'язкості
і параметра $y^+$ на достовірність CFD-розрахунків.

\tableofcontents
\newpage

% -------------------------------------------------------------------
\section{Вступ}

Турбулентність є однією з найскладніших проблем механіки рідини та газу
і водночас ключовим об'єктом сучасної обчислювальної гідродинаміки (CFD).
Практичні задачі --- від аеродинаміки літальних апаратів до теплообміну
в енергетичних установках --- майже завжди відбуваються при високих числах
Рейнольдса, коли потік стає хаотичним і тривимірним. Точне відтворення
таких процесів вимагає узгодженого поєднання фізичної моделі турбулентності,
математичного замикання рівнянь Нав'є--Стокса та чисельних методів
дискретизації.

% -------------------------------------------------------------------
\section{Фізична природа турбулентності та класичні теорії}

Турбулентність --- це складна, просторово-часова тривимірна нестаціонарна
форма руху рідини або газу, що характеризується вираженою хаотичністю та
високою інтенсивністю дисипації механічної енергії. На відміну від
ламінарного режиму, де шари середовища переміщуються впорядковано,
турбулентна течія супроводжується інтенсивним перемішуванням, пульсаціями
швидкості, тиску й температури, а також безперервним утворенням і руйнуванням
вихорів різного масштабу.

\subsection{Фундаментальні властивості}

\textbf{Хаотичність та непередбачуваність.} Локальні гідродинамічні параметри
демонструють випадкові коливання у часі та просторі, що унеможливлює точний
довгостроковий прогноз миттєвого стану течії і вимагає статистичних та
ймовірнісних методів опису.

\textbf{Просторова тривимірність.} Турбулентні пульсації швидкості завжди
тривимірні, навіть якщо осереднена течія є одновимірною чи двовимірною.
Це зумовлено механізмом розтягування вихрових трубок (\emph{vortex stretching})
--- нелінійним процесом, за якого кутова швидкість обертання вихорів зростає
при їхньому розтягуванні вздовж ліній середнього зсуву.

\textbf{Дисипативність.} Кінетична енергія, що надходить у систему від
зовнішніх джерел або середнього руху, безперервно переноситься вздовж
спектра вихорів і перетворюється на внутрішню теплову енергію під дією
молекулярного в'язкого тертя.

\subsection{Енергетичний каскад Річардсона--Колмогорова}

Згідно з цією концепцією, турбулентний потік складається з вихорів широкого
спектра розмірів. Енергія середньої течії генерує найбільші вихори з
характерним лінійним масштабом~$L$, визначеним геометрією задачі. У
інерційному інтервалі енергія передається «зверху вниз» практично без втрат
на в'язке тертя, доки розмір вихорів не досягає масштабу Колмогорова~$\eta$,
на якому домінують сили молекулярної в'язкості:
\begin{equation}
  \eta = \left( \frac{\nu^3}{\varepsilon} \right)^{1/4},
  \label{eq:kolmogorov}
\end{equation}
де $\varepsilon$ --- швидкість дисипації турбулентної кінетичної енергії.

\subsection{Число Рейнольдса}

Основним критерієм стійкості ламінарної течії та її переходу у турбулентний
стан виступає безрозмірне число Рейнольдса:
\begin{equation}
  \Rey = \frac{\rho U L}{\mu} = \frac{U L}{\nu},
  \label{eq:reynolds}
\end{equation}
де $\rho$ --- густина, $U$ --- характерна швидкість, $L$ --- характерний
масштаб, $\mu$ --- динамічна в'язкість, $\nu = \mu/\rho$ --- кінематична
в'язкість. При малих $\Rey$ молекулярна в'язкість демпфує збурення; при
перевищенні критичного $\Rey_{cr}$ сили інерції переважають, і течія
переходить у турбулентний режим.

% -------------------------------------------------------------------
\section{Пряме чисельне моделювання (DNS)}

Найбільш безкомпромісним методом є пряме чисельне моделювання (Direct
Numerical Simulation, DNS) --- безпосереднє розв'язання тривимірних
нестаціонарних рівнянь Нав'є--Стокса без емпіричних моделей турбулентності.
Це вимагає дискретизації процесів від масштабу~$L$ до масштабу~$\eta$.

Оскільки $\varepsilon \sim U^3/L$, з \eqref{eq:kolmogorov} випливає
\begin{equation}
  \eta \sim L\,\Rey^{-3/4}.
  \label{eq:eta-re}
\end{equation}
Для адекватного розв'язання крок сітки $\Delta x$ має бути сумірним з~$\eta$
($\Delta x \sim \eta$). Кількість вузлів в одному напрямку $n \sim L/\eta
\sim \Rey^{3/4}$, а загальна кількість просторових вузлів
\begin{equation}
  N = n^3 \sim \Rey^{9/4}.
  \label{eq:nodes}
\end{equation}

За умовою Куранта--Фрідріхса--Леві (CFL) часовий крок $\Delta t \sim \eta/U$,
тому кількість кроків за характерний час $T \sim L/U$ дорівнює
$N_t \sim L/\eta \sim \Rey^{3/4}$. Сумарна обчислювальна складність
\begin{equation}
  N_f = N \times N_t \propto \Rey^{9/4} \times \Rey^{3/4} = \Rey^3.
  \label{eq:complexity}
\end{equation}

\begin{table}[H]
  \caption{Теоретичні вимоги до ресурсів DNS залежно від числа Рейнольдса}
  \label{tab:dns}
  \centering
  \renewcommand{\arraystretch}{1.2}
  \begin{tabular}{|r|r|r|r|r|}
    \hline
    $\Rey$ & $n$ & $N$ & $N_t$ & $N_f$ (умовні флопи) \\
    \hline
    $10^2$   & 32    & $3{,}2 \times 10^4$   & 32    & $1{,}0 \times 10^6$ \\
    \hline
    $10^3$   & 178   & $5{,}6 \times 10^6$   & 178   & $1{,}0 \times 10^9$ \\
    \hline
    $10^4$   & 1000  & $1{,}0 \times 10^9$   & 1000  & $1{,}0 \times 10^{12}$ \\
    \hline
    $10^5$   & 5623  & $1{,}8 \times 10^{11}$ & 5623  & $1{,}0 \times 10^{15}$ \\
    \hline
    $10^6$   & 31622 & $3{,}2 \times 10^{13}$ & 31622 & $1{,}0 \times 10^{18}$ \\
    \hline
  \end{tabular}
\end{table}

Наведені оцінки показують, що DNS для інженерних задач ($\Rey \sim 10^6$--$10^8$)
на сучасному етапі є практично неможливим.

% -------------------------------------------------------------------
\section{Осереднення Рейнольдса та проблема замикання}

Для подолання обмежень DNS у прикладних задачах застосовують статистичний
опис на основі осереднення Рейнольдса. Миттєве значення гідродинамічної
змінної розкладається на осереднену та пульсаційну складові:
\begin{equation}
  u_i(\mathbf{x}, t) = \bar{u}_i(\mathbf{x}) + u_i'(\mathbf{x}, t),
  \qquad
  p(\mathbf{x}, t) = \bar{p}(\mathbf{x}) + p'(\mathbf{x}, t).
  \label{eq:reynolds-decomp}
\end{equation}
Осереднення за інтервалом $T_0$:
\begin{equation}
  \bar{u}_i(\mathbf{x}) = \frac{1}{T_0} \int_{0}^{T_0} u_i(\mathbf{x}, t)\, dt.
  \label{eq:time-avg}
\end{equation}

Основні статистичні тотожності Рейнольдса:
\begin{equation}
  \overline{u_i'} = 0,\quad
  \overline{\bar{u}_i} = \bar{u}_i,\quad
  \overline{\frac{\partial u_i}{\partial x_j}} = \frac{\partial \bar{u}_i}{\partial x_j},\quad
  \overline{u_i u_j} = \bar{u}_i \bar{u}_j + \overline{u_i' u_j'}.
  \label{eq:reynolds-identities}
\end{equation}

\subsection{Рівняння RANS}

Миттєві рівняння неперервності та збереження імпульсу для нестислої рідини:
\begin{equation}
  \frac{\partial u_i}{\partial x_i} = 0,
  \label{eq:continuity}
\end{equation}
\begin{equation}
  \rho \left( \frac{\partial u_i}{\partial t} + u_j \frac{\partial u_i}{\partial x_j} \right)
  = -\frac{\partial p}{\partial x_i} + \mu \frac{\partial^2 u_i}{\partial x_j \partial x_j}.
  \label{eq:ns-instant}
\end{equation}

Після осереднення одержуємо рівняння RANS:
\begin{equation}
  \rho \left( \frac{\partial \bar{u}_i}{\partial t} + \bar{u}_j \frac{\partial \bar{u}_i}{\partial x_j} \right)
  = -\frac{\partial \bar{p}}{\partial x_i}
  + \frac{\partial}{\partial x_j} \left[
      \mu \left( \frac{\partial \bar{u}_i}{\partial x_j} + \frac{\partial \bar{u}_j}{\partial x_i} \right)
      - \rho \overline{u_i' u_j'}
    \right].
  \label{eq:rans}
\end{equation}

Тензор $-\rho \overline{u_i' u_j'}$ називають тензором напружень Рейнольдса.
Він містить 6 незалежних компонент, тоді як система має лише 4 рівняння
(3 імпульсу + неперервність). Це --- \emph{проблема замикання} рівнянь
Нав'є--Стокса.

\subsection{Гіпотеза Буссінеска}

Найпоширеніше замикання --- напівкласична гіпотеза Буссінеска про вихрову
в'язкість:
\begin{equation}
  -\rho \overline{u_i' u_j'} = \mu_t \left( \frac{\partial \bar{u}_i}{\partial x_j}
    + \frac{\partial \bar{u}_j}{\partial x_i}
    - \frac{2}{3} \frac{\partial \bar{u}_k}{\partial x_k} \delta_{ij} \right)
    - \frac{2}{3} \rho k \delta_{ij},
  \label{eq:boussinesq}
\end{equation}
де $\mu_t$ --- динамічна турбулентна в'язкість,
$k = \tfrac{1}{2} \overline{u_i' u_i'}$ --- кінетична енергія турбулентності,
$\delta_{ij}$ --- символ Кронекера. Обмеження гіпотези --- припущення про
ізотропність $\mu_t$, що знижує точність у сильно закручених і анізотропних
течіях.

% -------------------------------------------------------------------
\section{Класифікація моделей турбулентності}

\subsection{Алгебраїчні моделі (zero-equation)}

Алгебраїчні моделі обчислюють $\mu_t$ без додаткових диференціальних
рівнянь переносу. За гіпотезою Прандтля для зсувного пристінкового шару:
\begin{equation}
  \nu_t = l_m^2 \left| \frac{\partial \bar{u}}{\partial y} \right|,
  \label{eq:prandtl}
\end{equation}
де $l_m$ --- довжина шляху перемішування (наприклад, $l_m = \kappa y$ біля
стінки, $\kappa \approx 0{,}41$).

Переваги: простота, мінімальні витрати, висока чисельна стійкість.
Недоліки: відсутність «пам'яті» турбулентності; складне емпіричне
налаштування $l_m$ для тривимірних течій.

\subsection{Двопараметричні RANS-моделі}

\subsubsection{Модель $k$--$\varepsilon$ (Launder \& Spalding, 1974)}

Замикаючі змінні --- $k$ та швидкість дисипації $\varepsilon$:
\begin{align}
  \frac{\partial (\rho k)}{\partial t} + \frac{\partial (\rho \bar{u}_i k)}{\partial x_i}
  &= \frac{\partial}{\partial x_j} \left[ \left( \mu + \frac{\mu_t}{\sigma_k} \right) \frac{\partial k}{\partial x_j} \right]
    + G_k - \rho \varepsilon, \label{eq:k-eps-k}\\
  \frac{\partial (\rho \varepsilon)}{\partial t} + \frac{\partial (\rho \bar{u}_i \varepsilon)}{\partial x_i}
  &= \frac{\partial}{\partial x_j} \left[ \left( \mu + \frac{\mu_t}{\sigma_\varepsilon} \right) \frac{\partial \varepsilon}{\partial x_j} \right]
    + C_{1\varepsilon} \frac{\varepsilon}{k} G_k - C_{2\varepsilon} \rho \frac{\varepsilon^2}{k}, \label{eq:k-eps-eps}
\end{align}
де $G_k = \mu_t S^2$, $S = \sqrt{2 S_{ij} S_{ij}}$, а
\begin{equation}
  \mu_t = \rho C_\mu \frac{k^2}{\varepsilon}.
  \label{eq:mut-keps}
\end{equation}
Стандартні константи: $C_\mu = 0{,}09$, $C_{1\varepsilon} = 1{,}44$,
$C_{2\varepsilon} = 1{,}92$, $\sigma_k = 1{,}0$, $\sigma_\varepsilon = 1{,}3$.
Модель ефективна для вільних зсувних течій, але слабка в пристінкових
шарах із позитивним градієнтом тиску (відрив потоку).

\subsubsection{Модель $k$--$\omega$ (Wilcox, 1988/2006)}

Замість $\varepsilon$ використовують питому швидкість дисипації
$\omega = \varepsilon/(\beta^* k)$:
\begin{align}
  \frac{\partial (\rho k)}{\partial t} + \frac{\partial (\rho \bar{u}_j k)}{\partial x_j}
  &= \frac{\partial}{\partial x_j} \left[ \left( \mu + \sigma_k \mu_t \right) \frac{\partial k}{\partial x_j} \right]
    + P_k - \beta^* \rho \omega k, \label{eq:k-omega-k}\\
  \frac{\partial (\rho \omega)}{\partial t} + \frac{\partial (\rho \bar{u}_j \omega)}{\partial x_j}
  &= \frac{\partial}{\partial x_j} \left[ \left( \mu + \sigma_\omega \mu_t \right) \frac{\partial \omega}{\partial x_j} \right]
    + \gamma \frac{\omega}{k} P_k - \beta \rho \omega^2, \label{eq:k-omega-omega}
\end{align}
$\mu_t = \rho k / \omega$. Переваги: інтегрування до стінки ($y \to 0$),
точність у граничному шарі. Недолік: чутливість до граничних умов $\omega$
у вільному потоці.

\subsubsection{Модель $k$--$\omega$ SST (Menter, 1994)}

Модель SST поєднує $k$--$\omega$ біля стінок і $k$--$\varepsilon$ у
вільному потоці через функцію змішування $F_1$ та член перехресної дифузії:
\begin{equation}
  CD_{k\omega} = \max \left( 2 \rho \sigma_{\omega 2} \frac{1}{\omega}
    \frac{\partial k}{\partial x_i} \frac{\partial \omega}{\partial x_i},\, 10^{-10} \right).
  \label{eq:cd-kw}
\end{equation}
\begin{align}
  \frac{\partial (\rho k)}{\partial t} + \frac{\partial (\rho \bar{u}_j k)}{\partial x_j}
  &= \frac{\partial}{\partial x_j} \left[ \left( \mu + \sigma_k \mu_t \right) \frac{\partial k}{\partial x_j} \right]
    + \tilde{P}_k - \beta^* \rho \omega k, \label{eq:sst-k}\\
  \frac{\partial (\rho \omega)}{\partial t} + \frac{\partial (\rho \bar{u}_j \omega)}{\partial x_j}
  &= \frac{\partial}{\partial x_j} \left[ \left( \mu + \sigma_\omega \mu_t \right) \frac{\partial \omega}{\partial x_j} \right]
    + P_\omega - \beta \rho \omega^2 \notag\\
  &\quad + 2(1 - F_1) \rho \sigma_{\omega 2} \frac{1}{\omega}
    \frac{\partial k}{\partial x_j} \frac{\partial \omega}{\partial x_j}. \label{eq:sst-omega}
\end{align}

Функція змішування:
\begin{equation}
  F_1 = \tanh \left( \arg_1^4 \right),
  \label{eq:f1}
\end{equation}
\begin{equation}
  \arg_1 = \min \left[
    \max \left( \frac{\sqrt{k}}{\beta^* \omega d},\; \frac{500\nu}{d^2 \omega} \right),\;
    \frac{4\rho\sigma_{\omega 2} k}{CD_{k\omega}\, d^2}
  \right],
  \label{eq:arg1}
\end{equation}
де $d$ --- відстань до найближчої стінки, $\beta^* = 0{,}09$. Константи моделі інтерполюються:
$\phi = F_1 \phi_1 + (1 - F_1)\phi_2$.

Обмеження в'язкості за Бредшоу:
\begin{equation}
  \mu_t = \frac{\rho a_1 k}{\max \left( a_1 \omega,\, S F_2 \right)},
  \qquad a_1 = 0{,}31,
  \label{eq:bradshaw}
\end{equation}
\begin{equation}
  F_2 = \tanh \left( \arg_2^2 \right),
  \quad
  \arg_2 = \max \left( \frac{2\sqrt{k}}{\beta^* \omega d},\; \frac{500 \nu}{d^2 \omega} \right).
  \label{eq:f2}
\end{equation}

SST є фактичним стандартом у аеродинамічному проектуванні, зокрема для
відривних течій (backward-facing step), де $k$--$\varepsilon$ занижує
відривну зону.

\subsection{Метод великих вихорів (LES)}

LES застосовує просторовий фільтр із характерним розміром~$\Delta$ (зазвичай
крок сітки). Фільтрована швидкість:
\begin{equation}
  \bar{u}_i(\mathbf{x}, t) = \int_{\Omega} G(\mathbf{x}, \mathbf{y}; \Delta)\, u_i(\mathbf{y}, t)\, d\mathbf{y}.
  \label{eq:filter}
\end{equation}

Фільтровані рівняння Нав'є--Стокса:
\begin{align}
  \frac{\partial \bar{u}_i}{\partial x_i} &= 0, \label{eq:les-cont}\\
  \frac{\partial \bar{u}_i}{\partial t} + \frac{\partial (\bar{u}_i \bar{u}_j)}{\partial x_j}
  &= -\frac{1}{\rho}\frac{\partial \bar{p}}{\partial x_i}
    + \nu \frac{\partial^2 \bar{u}_i}{\partial x_j \partial x_j}
    - \frac{\partial \tau_{ij}^{\mathrm{SGS}}}{\partial x_j}, \label{eq:les-momentum}
\end{align}
де $\tau_{ij}^{\mathrm{SGS}} = \overline{u_i u_j} - \bar{u}_i \bar{u}_j$ ---
тензор підсіткових напружень.

Замикання Смагоринського--Ліллі:
\begin{equation}
  \tau_{ij}^{\mathrm{SGS}} - \frac{1}{3}\tau_{kk}^{\mathrm{SGS}}\delta_{ij}
  = -2 \nu_{\mathrm{SGS}} \bar{S}_{ij},
  \qquad
  \nu_{\mathrm{SGS}} = (C_s \Delta)^2 |\bar{S}|,
  \label{eq:smagorinsky}
\end{equation}
$|\bar{S}| = \sqrt{2 \bar{S}_{ij} \bar{S}_{ij}}$, $C_s \approx 0{,}17$--$0{,}21$
(ізотропна турбулентність) або $0{,}065$--$0{,}1$ (зсувні течії). Динамічна
модель Germano--Lilly обчислює $C_s$ локально за тотожністю Germano.

% -------------------------------------------------------------------
\section{Вплив чисельних методів на моделювання турбулентності}

Точність CFD залежить не лише від моделі турбулентності, а й від властивостей
різницевих схем. Дослідження А.~Самарського та П.~Роуча показують жорсткий
зв'язок між локальною похибкою апроксимації та штучною дисипацією потоку.

\subsection{Схемна (чисельна) в'язкість}

Розглянемо одновимірне рівняння адвекції ($a > 0$):
\begin{equation}
  \frac{\partial u}{\partial t} + a \frac{\partial u}{\partial x} = 0.
  \label{eq:advection}
\end{equation}
Схема «вперед за часом, назад за простором» (upwind):
\begin{equation}
  \frac{u_j^{n+1} - u_j^n}{\Delta t} + a \frac{u_j^n - u_{j-1}^n}{\Delta x} = 0.
  \label{eq:upwind}
\end{equation}

Розклад у ряд Тейлора дає модифіковане рівняння (за Роучем):
\begin{equation}
  \frac{\partial u}{\partial t} + a \frac{\partial u}{\partial x}
  = \nu_{\mathrm{num}} \frac{\partial^2 u}{\partial x^2},
  \label{eq:modified}
\end{equation}
де
\begin{equation}
  \nu_{\mathrm{num}} = \frac{a \Delta x}{2} (1 - C),
  \qquad C = \frac{a \Delta t}{\Delta x}
  \label{eq:num-visc}
\end{equation}
--- число Куранта.

При $C < 1$ маємо $\nu_{\mathrm{num}} > 0$: стійкість забезпечується
штучним згладжуванням. Якщо $\nu_{\mathrm{num}}$ стає сумірною з
фізичною турбулентною в'язкістю $\nu_t$, чисельний алгоритм пригнічує
дрібні вихори --- виникає передчасна «ламінаризація». Тому конвективні
члени дискретизують схемами 2-го та вищих порядків (QUICK, TVD, WENO).

\subsection{Параметр $y^+$ і пристінкове моделювання}

Безрозмірна відстань від стінки:
\begin{equation}
  y^+ = \frac{u_\tau y}{\nu},
  \qquad
  u_\tau = \sqrt{\tau_w / \rho},
  \label{eq:yplus}
\end{equation}
де $y$ --- відстань до центру першої комірки, $\tau_w$ --- напруження тертя.

\begin{table}[H]
  \caption{Зони граничного шару та вимоги до сітки й моделей}
  \label{tab:yplus}
  \centering
  \small
  \renewcommand{\arraystretch}{1.15}
  \begin{tabular}{|l|c|p{4.5cm}|p{3.8cm}|}
    \hline
    Зона & $y^+$ & Особливості & Рекомендовані моделі \\
    \hline
    В'язкий підшар & $< 5$ &
      Домінує молекулярна в'язкість &
      $k$--$\omega$ SST, Wilcox $k$--$\omega$ \\
    \hline
    Буферний шар & $5$--$30$ &
      Суміш в'язких і турбулентних ефектів; уникати 1-го вузла тут &
      Універсальні wall functions \\
    \hline
    Логарифмічний шар & $30$--$500$ &
      Турбулентне перемішування, log-law &
      Стандартна $k$--$\varepsilon$ \\
    \hline
  \end{tabular}
\end{table}

Невідповідність сітки моделі (наприклад, $y^+ = 100$ для low-Re підходу)
приводить до грубих похибок у силах тертя та теплових потоках.

% -------------------------------------------------------------------
\section{Висновки}

\begin{enumerate}
  \item DNS --- єдиний безмодельний точний підхід, але обмежений низькими
        $\Rey$ через складність $N_f \propto \Rey^3$.
  \item RANS залишається основним інструментом для пристінкових інженерних
        течій. Найуніверсальнішою двопараметричною моделлю є $k$--$\omega$ SST,
        що поєднує переваги $k$--$\omega$ біля стінок і $k$--$\varepsilon$
        у вільному потоці.
  \item LES займає проміжне положення між RANS і DNS, дозволяючи моделювати
        нестаціонарну тривимірну динаміку великих вихорів, але вимагає
        значних обчислювальних ресурсів, особливо в пристінкових шарах.
  \item Фізично достовірний CFD-розрахунок потребує комплексного підходу:
        узгодженого вибору моделі турбулентності, контролю $y^+$ та схем
        підвищеного порядку для мінімізації схемної в'язкості.
\end{enumerate}

\begin{table}[H]
  \caption{Порівняльний аналіз підходів до моделювання турбулентності}
  \label{tab:compare}
  \centering
  \small
  \renewcommand{\arraystretch}{1.2}
  \begin{tabular}{|p{3.1cm}|c|c|p{4.2cm}|}
    \hline
    Клас моделі & Точність & Вартість & Обмеження \\
    \hline
    Алгебраїчні (zero-eq) & Дуже низька & Надзвичайно низька &
      Немає history effects \\
    \hline
    RANS ($k$--$\varepsilon$, $k$--$\omega$) & Середня & Помірна &
      Ізотропність $\mu_t$; чутливість до $y^+$ \\
    \hline
    $k$--$\omega$ SST & Висока (RANS) & Помірна &
      Не розв'язує нестаціонарні великі вихори \\
    \hline
    LES & Дуже висока & Висока &
      Високі вимоги до сітки та потужності \\
    \hline
    DNS & Абсолютна & Гранично висока &
      Лише низькі $\Rey$, прості геометрії \\
    \hline
  \end{tabular}
\end{table}

% -------------------------------------------------------------------
\section*{Список використаних джерел}
\addcontentsline{toc}{section}{Список використаних джерел}

\begin{enumerate}
  \item Pope S.~B.
  \emph{Turbulent Flows}.
  --- Cambridge University Press, 2000. --- 771~p.
  \item Wilcox D.~C.
  \emph{Turbulence Modeling for CFD}.
  --- DCW Industries, 2006. --- 522~p.
  \item Launder B.~E., Spalding D.~B.
  The numerical computation of turbulent flows //
  \emph{Computer Methods in Applied Mechanics and Engineering}. ---
  1974. --- Vol.~3. --- P.~269--289.
  \item Menter F.~R.
  Two-equation eddy-viscosity turbulence models for engineering applications //
  \emph{AIAA Journal}. --- 1994. --- Vol.~32, No.~8. --- P.~1598--1605.
  \item Smagorinsky J.
  General circulation experiments with the primitive equations //
  \emph{Monthly Weather Review}. --- 1963. --- Vol.~91. --- P.~99--164.
  \item Роуч П.
  \emph{Вычислительная гидродинамика}.
  --- М.: Мир, 1980. --- 616~с.
  \item Самарский А.~А.
  \emph{Теория разностных схем}.
  --- М.: Наука, 1989. --- 616~с.
  \item Андерсон Д., Таннехилл Дж., Плетчер Р.
  \emph{Вычислительная гидродинамика и теплообмен: в 2 т.}
  --- М.: Мир, 1990.
\end{enumerate}

\end{document}
